\documentclass[12pt]{article}
\usepackage[T1]{fontenc}
\usepackage[utf8]{inputenc}
\usepackage{lmodern}
\usepackage{textcomp}
\usepackage{enumitem}
\usepackage{geometry}
\geometry{margin=1in}
\begin{document}
\begin{center}
Answer Key - History - Graduate Level Assessment 9 \
\small (Answers are ordered exactly as in the question paper.)
\end{center}

\section*{Multiple Choice}
\begin{enumerate}[label=\arabic*.]
\item \textbf{Answer:} A
\\ \textit{Options:} A) China.; B) Africa.; C) Australia.; D) North America.; E) South America.
\\ \textit{Note:} The correct answer is China because archaeological evidence indicates that the earliest known production of bronze artifacts dates back to around 3000 BCE in the region of the Yellow River, marking a significant advancement in metallurgy during the Shang Dynasty.
\item \textbf{Answer:} A
\\ \textit{Options:} A) focus on creating intricate geometric patterns.; B) construction of structures exclusively on hilltops.; C) employment of advanced hydraulic systems.; D) incorporation of wood, earth, and stone.; E) reliance on natural materials such as leaves and vines.
\\ \textit{Note:} The correct answer highlights that Inca architecture is characterized by its sophisticated use of geometric patterns, which reflect their advanced engineering skills and aesthetic principles, as seen in their temples and agricultural terraces.
\item \textbf{Answer:} C
\\ \textit{Options:} A) a hieroglyphic written language and calendar.; B) stratified societies ruled by kings.; C) a sophisticated bronze producing industry.; D) construction of large-scale pyramids.; E) usage of gold currency.
\\ \textit{Note:} The correct answer is a sophisticated bronze producing industry because the Maya civilization primarily utilized stone tools and did not develop metallurgy to the extent of producing bronze before 1900 B.P.
\item \textbf{Answer:} D
\\ \textit{Options:} A) severe drought and a surge in warfare; B) increased rainfall and population migration to the Maya ceremonial centers; C) severe drought and increased storage of surplus food from agriculture; D) severe drought and a decline in construction at Maya ceremonial centers; E) increased warfare and the abandonment of the entire region
\\ \textit{Note:} The correct answer reflects the correlation between periods of severe drought, indicated by variations in oxygen isotopes, and the subsequent decline in construction activity at Maya ceremonial centers, highlighting the impact of environmental changes on societal developments.
\item \textbf{Answer:} A
\\ \textit{Options:} A) Europe; B) Africa; C) Central America; D) Asia; E) None of the above
\\ \textit{Note:} Europe features the earliest evidence of etching into a rock face, as demonstrated by archaeological findings in sites such as the Gravettian culture in the Upper Paleolithic period, where intricate carvings on rocks and cave walls have been dated to around 30,000 to 40,000 years ago.
\end{enumerate}

\section*{True / False}
\begin{enumerate}[label=\arabic*.]
\setcounter{enumi}{5}
\item \textbf{Answer:} True
\\ \textit{Options:} A) True; B) False
\\ \textit{Note:} According to the passage: the dissolution and reconstitution of three-dimensional form
\item \textbf{Answer:} False
\\ \textit{Options:} A) True; B) False
\\ \textit{Note:} The correct answer is: hunting of various animals
\item \textbf{Answer:} False
\\ \textit{Options:} A) True; B) False
\\ \textit{Note:} The correct answer is: Jules J. Janssens
\item \textbf{Answer:} True
\\ \textit{Options:} A) True; B) False
\\ \textit{Note:} According to the passage: 100,000
\item \textbf{Answer:} True
\\ \textit{Options:} A) True; B) False
\\ \textit{Note:} According to the passage: The Premier League is particularly popular in Asia, where it is the most widely distributed sports programme
\end{enumerate}

\section*{Long Form Response}
\begin{enumerate}[label=\arabic*.]
\setcounter{enumi}{10}
\item \textbf{Answer:} where content protection code instructs the player to refuse commands to skip a specific part
\\ \textit{Note:} Expected answer: where content protection code instructs the player to refuse commands to skip a specific part
\item \textbf{Answer:} independent republics
\\ \textit{Note:} Expected answer: independent republics
\end{enumerate}

\vfill
\noindent\textit{End of Answer Key}
\end{document}